একজন শিক্ষার্থী ঘরে বসে বিদেশের পাঠে অংশ নেয়, গ্রামের রোগী দূরের চিকিৎসকের পরামর্শ পান, আর স্থানীয় উদ্যোক্তা অনলাইনে পণ্য বিক্রি করেন। একই তথ্য ও যোগাযোগ অবকাঠামো বিভিন্ন কাজে সময় ও দূরত্বের বাধা কমায়। এই অধ্যায়ে সে অবকাঠামোর ধারণা, সংশ্লিষ্ট প্রযুক্তি এবং দায়িত্বশীল ব্যবহারের প্রয়োজন ব্যাখ্যা করা হয়েছে।
\begin{learningbox}
এই অধ্যায় শেষে শিক্ষার্থী পারবে:
\begin{multicols}{2}
\begin{itemize}
\item ডেটা, তথ্য, IT ও ICT-এর সম্পর্ক এবং বিশ্বগ্রামের পাঁচ উপাদান ব্যাখ্যা করতে।
\item যোগাযোগ, কর্মসংস্থান, শিক্ষা, চিকিৎসা, গবেষণা, অফিস, বাসস্থান, ব্যবসা, সংবাদ, বিনোদন ও সাংস্কৃতিক বিনিময়ে ICT-এর ব্যবহার বুঝতে।
\item VR ও AR-এর পার্থক্য, প্রয়োজনীয় যন্ত্র এবং নিরাপদ প্রশিক্ষণের উপযোগিতা লিখতে।
\item AI, Expert System ও Robotics-এর কাজ এবং মানুষের তত্ত্বাবধানের প্রয়োজন ব্যাখ্যা করতে।
\item Cryosurgery, মহাকাশ গবেষণা, উৎপাদন ও প্রতিরক্ষায় ICT-এর ভূমিকা লিখতে।
\item Biometrics-এর বৈশিষ্ট্য ও পরিচয় যাচাইয়ের ধাপ শনাক্ত করতে।
\item Bioinformatics-এর তথ্য বিশ্লেষণ এবং Genetic Engineering-এর DNA পরিবর্তনের পার্থক্য বুঝতে।
\item Nanotechnology-এর একক, পদ্ধতি, ব্যবহার ও সীমাবদ্ধতা ব্যাখ্যা করতে।
\item গোপনীয়তা, স্বত্ব, plagiarism, cybercrime ও নিরাপদ ব্যবহারের নীতি প্রয়োগ করতে।
\item প্রযুক্তির সামাজিক ও অর্থনৈতিক প্রভাব এবং ব্যবহারের সুবিধা ও সীমাবদ্ধতা বিশ্লেষণ করতে।
\end{itemize}
\end{multicols}
\end{learningbox}
\textbf{পড়ার নিয়ম:} সংজ্ঞা জানবে, কাজের ধাপ বুঝবে, উদাহরণে প্রয়োগ করবে, তারপর সুবিধার সঙ্গে সীমাবদ্ধতা বিচার করবে। অতিরিক্ত আধুনিক উদাহরণ মূল ধারণা বোঝার সহায়ক; সেগুলোর পণ্যের নাম বা সংস্করণ মুখস্থ করা অধ্যায়ের মূল লক্ষ্য নয়।
